\documentclass[10pt]{article}
\usepackage[margin=20mm]{geometry}
\usepackage{amsmath,amssymb,mathtools,fontspec,unicode-math,xurl,hyperref}
\IfFontExistsTF{DejaVu Serif}{\setmainfont{DejaVu Serif}}{\setmainfont{Latin Modern Roman}}
\setmathfont{latinmodern-math.otf}
\setlength{\parindent}{0pt}\setlength{\parskip}{6pt}\sloppy
\title{Technical Verification of
Toric Stratifications of Character Varieties}
\author{AI-assisted technical review}
\date{}
\begin{document}
\maketitle
Abstract: The acquired publisher version is reviewed through its stratification, surface-pairing and Springer-action arguments. One local computational defect is established: the closed formula in Proposition 2.21 reconstructs incorrect compact-support Hodge numbers even for a two-dimensional torus. The primitive-block inversion is valid, and the correct geometric sum restores that proposition. The preceding sum in the printed proof is already correct; no dependence of the main curious Lefschetz argument on the erroneous closed form was found. The supported assessment is Minor. Integral lattice saturation, affine shifts, the Springer sign convention and flatness on the weight-graded cohomology were checked at their stated interfaces. This initial review does not constitute independent acceptance; scientific adjudication remains pending.\par
\section*{Identity, exact version and scope}
The paper is Anton Mellit, Toric stratifications of character varieties, Publications Mathématiques de l’IHÉS 142 (2025), 153–240, DOI 10.1007/s10240-025-00158-0 (https://doi.org/10.1007/s10240-025-00158-0) [paper]. This report concerns the acquired 88-page publisher version of record, published online on 25 September 2025. Physical PDF page \(p\) has printed number \(p+152\). The separately acquired arXiv:1905.10685v1 source, dated 25 May 2019 and titled Cell decompositions of character varieties, is linked from the author’s publication entry. It is an earlier related version; neither equality with the publisher text nor identity as its typesetting source is asserted.\par
The review reconstructs the main proof dependencies and selected calculations, with actual source-image inspection of the formulas and diagrams used as evidence. Complete text reading does not mean that every external theorem was independently reproved or that every source page was visually inspected. The coverage and source-page records retain those distinctions. No separate supplement, computational release or correction was identified in the bounded publisher, author-page and arXiv checks. Their absence is not treated as a mathematical defect.\par
\section*{Definitions and main proof dependencies}
The operative genericity condition is Definition 4.9, PDF p. 40: \(\prod_i\det C_i=1\), and for every \(1\le r<n\), every choice of \(r\) eigenvalues from each \(C_i\) has total product different from one. A common invariant proper subspace would violate this condition by taking determinants on that subspace. It therefore excludes reducible representations, rather than merely imposing a condition on the full determinant. The introductory omission of \(r=1\), and the repeated subscript in the determinant product, are resolved by the body definition and are not additional Problems.\par
The smoothness and scalar-stabilizer interface in Theorem 4.10 agrees with the explicitly consulted author version of Hausel–Letellier–Rodriguez-Villegas [hlrv], Definition 2.1.1, Proposition 2.1.4 and Theorem 2.1.5, PDF pp. 16–18. Nonemptiness is kept separate from smoothness of a possibly empty variety. At \(k=0\), the genericity condition excludes \(n>1\); when \(n=1\), the resulting space is \((\mathbb C^*)^{2g}\). Thus Theorem 8.13 does not assert an unrestricted result for all untwisted higher-rank character varieties.\par
Equation reference: eq:middle. For a regular semisimple last puncture, set \(r=\binom n2\) and \(d=\dim_{\mathbb C}X\). Theorem 1.2 gives a rank-\(r\) vector bundle over \(X\), with strata \[(\mathbb C^*)^{d-2i}\times\mathbb C^{i+r}.\] The center of the compact-support curious Lefschetz symmetry on each stratum is \[
 (d-2i)+2(i+r)=d+2r.\] The affine-bundle isomorphism shifts cohomological degree and weight by \(2r\), recovering center \(d\) for \(X\). Consequently Proposition 2.19 is applied with a common center on all strata. The long exact compact-support sequence is compatible with cup product, and strictness of mixed Hodge structures makes its weight-graded rows exact. The five-lemma argument then preserves the asserted isomorphisms.\par
\section*{A local defect in Hodge-number reconstruction}
Problem 1 (The geometric sum in Proposition 2.21 has the wrong ratio).  In Proposition 2.21, PDF p. 21 (printed p. 173), both the statement and the last closed-form step in its proof use a geometric sum with ratio \(uv\). A primitive Lefschetz step raises both the cohomological degree and the Hodge index by two, so the ratio must be \(u^2v^2\).\par
\subsection*{Claim and hypotheses}
Equation reference: eq:inversion, eq:printed. Write \(N=\cup\omega\), assume that \(X\) is Tate, \(\omega\in F^2H^2(X)\), and \(N\) satisfies curious Lefschetz with middle weight \(2d\). The proposition defines \[g_a(u)=\sum_k u^k\dim\ker\bigl(N^a:H_c^k(X)\longrightarrow
 H_c^{k+2a}(X)\bigr)\quad(a>0),\qquad g_a=0\quad(a\le0),\] and \[
 f_i(u)=(1+u^2)g_{i+1}(u)-g_i(u)-u^2g_{i+2}(u).\] Its claimed output is \[
 \operatorname{MH}_X(u,v)=\sum_{i\ge0}v^{d-i}f_i(u)
       \frac{(uv)^{i+1}-1}{uv-1},
 \qquad
 \operatorname{MH}_X(u,v)=\sum_{k,j}u^kv^j\dim H_c^k(X)^{j,j}.\] The quotient is a finite polynomial sum, so the issue is unrelated to substitution at \(uv=1\).\par
\subsection*{A counterexample satisfying every hypothesis}
Take \(X=(\mathbb C^*)^2\) and \(\omega=d\log x\wedge d\log y\). This smooth variety is Tate, and the nonzero class of \(\omega\) lies in \(F^2H^2(X)\). Its compact-support groups are \[H_c^2(X)=\mathbb C(0),\qquad
 H_c^3(X)=\mathbb C(-1)^{\oplus2},\qquad
 H_c^4(X)=\mathbb C(-2).\] Poincaré duality shows that \(N:H_c^2\to H_c^4\) is an isomorphism; all other positive-degree maps are zero for degree reasons. Thus curious Lefschetz holds with middle weight \(2\), so \(d=1\). This example is also a rank-one genus-one character variety.\par
Equation reference: eq:wrongoutput, eq:trueoutput. The actual kernel polynomials are \[g_1=2u^3+u^4,\qquad
 g_a=u^2+2u^3+u^4\quad(a\ge2).\] Equation [eq:inversion] (\#eq:inversion) correctly yields \(f_0=2u^3,\ f_1=u^2\), and \(f_i=0\) for \(i\ge2\). However, substitution in [eq:printed] (\#eq:printed) gives \[
 \operatorname{MH}_{\rm printed}=2u^3v+u^2(1+uv)=u^2+3u^3v,\] whereas the actual answer is \[
 \operatorname{MH}_X=u^2+2u^3v+u^4v^2.\] The printed output assigns dimension three to \(H_c^3\) of type \((1,1)\), and zero to the nonzero top group \(H_c^4\). These are incorrect numerical invariants, not merely alternative notations. The exact calculation and its input hypotheses are saved in `verification/0137/root\_prop221\_exact.json`.\par
\subsection*{General correction and its verification}
Equation reference: eq:correct. Let \(F_m(u)\) be the graded polynomial of primitive vectors in blocks of length \(m+1\). With \(\Delta g_a=g_a-g_{a-1}\), direct inspection of a nilpotent block gives \[\Delta g_a=\sum_{m\ge a-1}u^{2(m-a+1)}F_m,\qquad
 \Delta g_{i+1}-u^2\Delta g_{i+2}=F_i.\] This proves [eq:inversion] (\#eq:inversion). A primitive vector with Hodge index \(d-i\) contributes its successive powers under \(N\) at indices \(d-i,d-i+2,\ldots,d+i\). Therefore the corrected statement is \[
 \operatorname{MH}_X(u,v)
 =\sum_{i\ge0}v^{d-i}f_i(u)\sum_{j=0}^i(u^2v^2)^j
 =\sum_{i\ge0}v^{d-i}f_i(u)
       \frac{(uv)^{2(i+1)}-1}{(uv)^2-1}.\] The finite-sum interpretation applies also at \((uv)^2=1\). This derivation verifies the correction for every primitive block under the original assumptions, and hence for their direct sum. The saved exact symbolic checks additionally reproduce the torus counterexample and the kernel inversion on 35 finite blocks.\par
\subsection*{Counterevidence, impact and classification}
The strongest counterevidence is on the same page: the preceding double sum in the printed proof already contains \(v^{d-i+2j}u^{2j}f_i(u)\), which is correct. The primitive decomposition and kernel inversion also remain valid. The error is therefore consistent with a transcription mistake. No subsequent use of the incorrect closed form was found in the canonical text, and the main stratification and curious Lefschetz proofs use other propositions.\par
This report retains one Minor finding because the independently stated numerical reconstruction, when used as printed, produces the wrong Hodge numbers on a fully admissible input. It does not count unrelated typographical slips as findings. The supported consequence is confined to Proposition 2.21’s closed formula. Neither a failure of the main Lefschetz theorem nor a defect in its primitive-block mechanism is established. Equation [eq:correct] (\#eq:correct) restores the stated numerical reconstruction; author adoption and any prior discovery remain unverified.\par
\section*{Geometric and representation-theoretic checks}
The bounded checks of Sections 5–6 include the braid relations, the up/down/stay stratification, the local surface gluings and the identification of the logarithmic two-form with the intersection pairing on the capped surface, up to its fixed nonzero normalization. The overline in \(H_1(\overline\Sigma)\) in Corollary 6.9 is present in the actual image; its absence in extracted text is not a mathematical defect.\par
Two integral points in the quotient argument deserve explicit treatment. For Proposition 6.10, the boundary map \[H_1(\Sigma,A;\mathbb Z)\longrightarrow H_0(A;\mathbb Z)\] has image \(\ker(H_0(A;\mathbb Z)\to H_0(\Sigma;\mathbb Z))\), a saturated lattice. Its integral dual controls the stabilizer. This excludes a hidden finite stabilizer, which a dimension count alone would not exclude. For Proposition 6.13, the nonzero coordinate weights are differences of the torus coordinates attached to cycles of the braid permutation. Freeness makes the associated graph connected; a spanning tree supplies an integral character basis. Those coordinates give the asserted Zariski local trivialization. A general free torus representation would not justify choosing such a subset without this root-weight structure.\par
Theorem 6.14 has the common center \[\dim H_1(\overline\Sigma)+2|U_p|
   =2-n+|S_p|-c(\pi)+2|U_p|
   =2-n+\ell(\beta)-c(\pi).\] Here \(|S_p|+2|U_p|=\ell(\beta)\). The genericity condition rules out disconnected surfaces. The genus contributions and affine shifts in Section 7 were checked against this center and the dimension of the character variety. As a finite positive control, the rank-two four-puncture example has four connected walks at braid length six and one at length four. With the three choices at length four this gives seven strata and \[E_c(X;q)=(q-1)^2+6q=q^2+4q+1.\] This verifies the stated finite combinatorics and additive invariant, without asserting a new proof of the geometric stratification.\par
The descent of the two-form uses equality on the full action groupoid, which includes horizontality, rather than invariance under constant group elements alone. The torus splitting argument and strictness give its \(F^2\) property. The consulted author version of Guruprasad–Huebschmann–Jeffrey–Weinstein [ghjw], Theorems 7.1 and 8.3, supplies the nondegenerate parabolic pairing interface. Its invariant bilinear form has a normalization choice. For example, Mellit’s rank-one genus-one formula is \(2\,d\log a\wedge d\log b\); this corresponds to scaling the trace pairing. A nonzero scalar affects neither nondegeneracy nor curious Lefschetz. Literal identity of all conventions in the two published versions is not claimed.\par
\section*{The general-puncture passage}
Proposition 8.7 is checked on the full generic open set \(V\), including the repeated-eigenvalue fiber at one. After the affine cover, each relevant stratum factors over \(V\) as an affine space times a torus times \(V\). The saturated boundary lattice supplies the required splitting of the torus map. Compact-support pushforwards of the pieces are constant; successive extensions give locally constant cohomology sheaves. Strict weight filtrations retain the corresponding conclusion for the associated graded. Constancy of individual Hodge subspaces is not assumed.\par
For the Springer convention, the exact author version arXiv:1304.2642v2 of Achar–Henderson–Juteau–Riche [ahjr], Lemmas 4.3, 4.6 and 4.7, identifies the restriction/Grothendieck action on the top flag cohomology as the sign character. Its differently defined compact-support action must not be substituted. For \(\operatorname{GL}_n\), translation \(g\mapsto g-I\) identifies the relevant group neighborhood with an open subset of the Lie algebra, preserving stable flags and the regular-semisimple covering. This supplies the group-to-Lie-algebra interface.\par
Let \(N=\dim(G/B)\) in this paragraph. The sign summand in Proposition 8.5 is supported at the identity and is \(i_*\mathbb C[-2N]\); integration along the flag variety selects it. The zero-map assertion can be justified by \(\operatorname{Hom}(A,i_*V)=\operatorname{Hom}(i^*A,V)\). It does not require the false general principle that every derived morphism vanishing on all stalk cohomology groups is zero. Corollary 8.11 shifts both degree and weight by \(2N\), as required.\par
Section 8.5 needs flatness only after passing to \(\operatorname{Gr}^W_4H^2\). The displayed derivative is represented by \(2\sum_i d\log t_i\otimes c_1(L_i)\). Individual flag quotients may require a scalar-gauge choice. Intrinsically, the ratios \(R_i=L_i\otimes L_1^{-1}\) descend, and \(\sum_i d\log t_i=0\) on \(\det t=1\) gives \[\nabla[\omega]=
 2\sum_{i\ge2}d\log t_i\otimes c_1(R_i).\] Its values lie in \(W_2\). Thus the weight-graded cup operator is parallel and commutes with the relevant Weyl action. Projection to the sign part preserves Lefschetz isomorphisms, and the center shifts by \(2N\); the difference of fiber dimensions is exactly \(n^2-n=2N\). These checks support the stated general-puncture argument within the declared reference interfaces.\par
\section*{Conclusion and limits}
One local numerical reconstruction defect is supported, with a complete admissible example and a general algebraic correction. No substantive failure of Theorems 1.2, 1.5, 1.6, 7.5 or 8.13 was established in the inspected scope. The review retains mixed Hodge theory, Grothendieck–Springer theory and the cited smooth-quotient results as identified external inputs; it is not a recursive verification of all their foundations. The precise author versions consulted are recorded below. The bounded availability check does not establish that no later correction or additional source exists.\par
Dimension | Recorded state\par
Issue assessment | Minor: Proposition 2.21’s closed formula.\par
Repair verification | Equation [eq:correct] (\#eq:correct) follows for every primitive block under the original hypotheses.\par
Main proof impact | No dependency on the erroneous closed form identified; no main-theorem counterexample established.\par
Author adoption | Unknown; no author contact or adoption claim.\par
Scientific adjudication | Pending; human-review list empty.\par
Independent acceptance | Not performed by this initial review.\par
9 A. Mellit, Toric stratifications of character varieties, Publ. Math. IHÉS 142 (2025), 153–240. doi:10.1007/s10240-025-00158-0 (https://doi.org/10.1007/s10240-025-00158-0). Reviewed publisher version. T. Hausel, E. Letellier and F. Rodriguez-Villegas, Arithmetic harmonic analysis on character and quiver varieties, arXiv:0810.2076v3 (https://arxiv.org/abs/0810.2076v3), 16 September 2011. Exact author version consulted at the stated interface. K. Guruprasad, J. Huebschmann, L. Jeffrey and A. Weinstein, Group systems, groupoids, and moduli spaces of parabolic bundles, arXiv:dg-ga/9510006v1 (https://arxiv.org/abs/dg-ga/9510006v1). Exact author version consulted; publisher-version equivalence not asserted. P. N. Achar, A. Henderson, D. Juteau and S. Riche, Weyl group actions on the Springer sheaf, arXiv:1304.2642v2 (https://arxiv.org/abs/1304.2642v2), 3 October 2013. Exact accepted author version, PDF pp. 12–14.\par
\section*{Supporting checks}
The canonical PDF SHA256, with the two lines concatenated, is \[\texttt{b30acba27823b676237fb7bbcfc3528e}\] \[\texttt{51851bb68baa71f7b25b4914ae7d9a01}.\] The complete 64-character value is also retained programmatically in the source manifest and verification records. Own exact computations include the torus counterexample, 35 primitive blocks, 15 exterior Lefschetz matrices and 140 instances of [eq:middle] (\#eq:middle). Finite positive controls supplement the general arguments and do not certify unrestricted ranks. No author code or author-source compilation was used.\par
Their contributions are AI cross-checks, not human reviews or independent acceptance. No independent UI observation or backend equivalence is asserted.\par
\end{document}
